%%%%%%%%%%%%%%%%%%%%%%%%%%%%%%%%%%%%%%%%%%%%%%%%%%%%%%%%%%%%%%%%%%%%%%
% Lớp 11 | Xác suất | Biến cố giao và quy tắc nhân xác suất
% Dạng 1D9?1-4: Tính xác suất biến cố giao bằng sơ đồ hình cây
% Số câu: 2  (sinh tự động bởi tach_cau.py; từ nay thầy sửa trực tiếp file này)
%%%%%%%%%%%%%%%%%%%%%%%%%%%%%%%%%%%%%%%%%%%%%%%%%%%%%%%%%%%%%%%%%%%%%%

%%%%%%%%%%%%%%%% ĐÚNG/SAI %%%%%%%%%%%%%%%%

% nguồn: Ok_dungsai_lop11_Lan1
\begin{ex}%[1D9H1-4]%[ma:d538645e]
 Một xạ thủ bắn lần lượt $2$ viên đạn vào một  bia. Xác suất trúng đích của viên thứ nhất và thứ hai lần lượt là $0{,}9$ và $0{,}6$. Biết rằng kết quả các lần bắn là độc lập với nhau. Khi đó, xác suất của biến cố
 \choiceTF 
  {\lq\lq Cả 2 lần  bắn đều trúng đích\rq\rq\, là $0{,}04$}
  {\lq\lq Cả 2 lần bắn đều không trúng đích\rq\rq\, là $0{,}06$} 
  {\lq\lq Lần bắn thứ nhất trúng đích,  lần bắn thứ hai không trúng đích\rq\rq\, là $0{,}54$}
  {\lq\lq Lần bắn thứ nhất không trúng đích, lần bắn thứ hai  trúng đích\rq\rq\, là $0{,}38$}
  \loigiai{
  Gọi $A$ là biến cố \lq\lq Xạ thủ bắn viên  đạn trúng bia lần thứ nhất\rq\rq. Ta có $\mathrm{P}(A)=0{,}9$ và $\mathrm{P}(\overline{A})=0{,}1$.\\
  Gọi $B$ là biến cố \lq\lq Xạ thủ bắn  viên đạn trúng bia lần thứ hai\rq\rq. Ta có $\mathrm{P}(B)=0{,}6$ và $\mathrm{P}(\overline{B})=0{,}4$.\\
  Ta cũng có thể giải bài toán trên bằng cách sử dụng sơ đồ  hình cây như sau:
  \begin{center}
   \begin{tikzpicture}[font=\footnotesize,line join=round,line cap=round,>=stealth,scale=1,rotate=-90]
      \path (0,0) node (a){Xác suất}
      ++(2,4) node (b){$ A $ trúng bia  lần 1}
       (a)++(2,-4) node (c){$ \overline{A} $ không trúng bia lần 1};
     \path 
      (b)++(2,2) node  (b1){$ B $ trúng bia lần 2}
      (b)++(2,-2) node (b2){$ \overline{B} $ không trúng bia  lần 2};
     \path (c)++(2,2) node (c1){$ B $ trúng bia  lần 2}
      (c)++(2,-2) node (c2){$\overline{B} $ không  trúng bia lần 2};
      \path
       (b1)++(2,0) node[rectangle,draw=brown,text width=2cm,rounded
      corners=2mm,align=center](b11){Bắn trúng lần  1 và bắn trúng lần 2};
      \path
      (b2)++(2,0) node[rectangle,draw=brown,text  width=2cm,rounded
      corners=2mm,align=center](b21){Bắn trúng  lần 1 và không bắn trúng lần 2};
     \path (c1)++(2,0) node[rectangle,draw=brown,text 
      width=2cm,rounded 
       corners=2mm,align=center](c11){Không bắn trúng lần 1 và bắn trúng lần 2};
      \path (c2)++(2,0) node[rectangle,draw=brown,text
       width=2cm,rounded
      corners=2mm,align=center](c21){Không bắn trúng lần 1 và  không bắn trúng lần 2};
     \draw[-stealth,teal,outer  sep=0,inner
      sep=0](a.south)--(b.north) node[pos=0.5,right]{$0{,}9$}; 
      \draw[-stealth,teal,outer sep=0,inner
      sep=0](a.south)--(c.north) node[pos=0.5,left]{$0{,}1$}; 
      \draw[-stealth,teal,outer sep=0,inner
       sep=0](b.south)--(b1.north) node[pos=0.5,right]{$0{,}6$};
     \draw[-stealth,teal,outer  sep=0,inner
       sep=0](b.south)--(b2.north) node[pos=0.5,left]{$0{,}4$};
     \draw[-stealth,teal,outer  sep=0,inner
      sep=0](c.south)--(c1.north) node[pos=0.5,right]{$0{,}6$}; 
     \draw[-stealth,teal,outer  sep=0,inner
      sep=0](c.south)--(c2.north)  node[pos=0.5,left]{$0{,}4$};
     \draw[-stealth,teal,outer  sep=0,inner
      sep=0](b1.south)--(b11.north)  node[pos=0.5,left]{$0{,}54$};
      \draw[-stealth,teal,outer sep=0,inner
      sep=0](c1.south)--(c11.north)  node[pos=0.5,left]{$0{,}06$};
     \draw[-stealth,teal,outer  sep=0,inner
      sep=0](b2.south)--(b21.north) node[pos=0.5,left]{$0{,}36$}; 
     \draw[-stealth,teal,outer  sep=0,inner
      sep=0](c2.south)--(c21.north) node[pos=0.5,left]{$0{,}04$}; 
    \end{tikzpicture}
  \end{center}
  Theo sơ đồ trên  thì
  \begin{itemchoice}
   \itemch Xác suất cả 2 lần bắn đều  trúng đích là $0{,}54$.
   \itemch Xác suất cả  2 lần bắn đều không trúng đích là $0{,}04$.
   \itemch Xác suất lần  bắn thứ nhất trúng đích, lần bắn thứ hai không trúng đích là $0{,}36$.
   \itemch Xác suất lần bắn thứ nhất không trúng đích, lần bắn thứ hai trúng đích là  $0{,}06$.
  \end{itemchoice}
 } 
\end{ex}

% nguồn: Ok_dungsai_lop11_Lan1
\begin{ex}%[1D9H1-4]%[ma:eec147ca]
 Minh gieo $1$ hạt đậu và $1$ hạt ngô. Xác  suất nảy mầm của hạt đậu và hạt ngô lần lượt là $0{,}7$ và $0{,}6$. Biết rằng sự nảy mầm của hai hạt này là độc lập. Sử dụng sơ đồ hình cây, xác suất của các biến cố
  \choiceTF
  {\True \lq\lq  Cả $2$ hạt đều nảy mầm\rq\rq\, là $0{,}42$}
  {\True \lq\lq Cả $2$ hạt đều không nảy  mầm\rq\rq\, là $0{,}12$}
  {\lq\lq Hạt đậu nảy mầm, hạt ngô không  nảy mầm\rq\rq\, là $0{,}18$}
  {\lq\lq Hạt  đậu không nảy mầm, hạt ngô nảy mầm\rq\rq\, là $0{,}28$}
  \loigiai{
  Ta có sơ đồ hình cây  như sau:
  \begin{center}
   \begin{tikzpicture}[>=stealth,line cap=round,line join=round]
     \path 
      (0,0) coordinate (O) 
      (1,2) coordinate  (A)
      (1,-2) coordinate (B) 
       (3,2) coordinate (A1)
      (4,3)  coordinate (A2)
      (4,1)  coordinate (A2')
      (3,-2)  coordinate (B1)
      (4,-1) coordinate (B2) 
      (4,-3) coordinate (B2') 
       (7,4) coordinate (C) node[right]{Kết quả}
      (7,3) coordinate (C1) node[above right]{Hạt đậu nảy mầm,}node[below right]{Hạt ngô  nảy mầm}
      (7,1) coordinate (C2) node[above right]{Hạt đậu nảy mầm,}node[below right]{Hạt  ngô không nảy mầm}
      (7,-1) coordinate (C3) node[above right]{Hạt đậu không  nảy mầm,}node[below right]{Hạt ngô nảy mầm}
      (7,-3) coordinate  (C4) node[above right]{Hạt đậu không nảy mầm,}node[below right]{Hạt ngô không nảy mầm}
       (11.5,4) coordinate (D) node[right]{Xác suất}
      (11.5,3) coordinate (D) node[right]{$0{,}42$} 
      (11.5,1) coordinate (D)  node[right]{$0{,}28$}
       (11.5,-1) coordinate (D) node[right]{$0{,}18$}
      (11.5,-3) coordinate  (D) node[right]{$0{,}12$};
     \draw[->](O)--(A) node[above right]{Hạt đậu}node[below right]{nảy mầm}node[midway,left]{$0{,}7$}; 
     \draw[->](A1)--(A2) node[above  right]{Hạt ngô}node[below right]{nảy mầm}node[midway,left]{$0{,}6$};
      \draw[->](A1)--(A2') node[above right]{Hạt ngô}node[below right]{không nảy mầm}node[midway,left]{$0{,}4$};
     \draw[->](O)--(B) node[above right]{Hạt đậu}node[right]{không}node[below  right]{nảy mầm}node[midway,left]{$0{,}3$};
     \draw[->](B1)--(B2) node[above right]{Hạt ngô}node[below right]{nảy mầm}node[midway,left]{$0{,}6$}; 
      \draw[->](B1)--(B2') node[above right]{Hạt ngô}node[below right]{không nảy mầm}node[midway,left]{$0{,}4$};
    \end{tikzpicture}
  \end{center}
  \begin{itemchoice}
   \itemch Xác suất cả  $2$ hạt đều nảy mầm là $0{,}42$.
   \itemch Xác suất cả $2$ hạt đều không nảy  mầm là $0{,}12$.
   \itemch Xác suất hạt đậu nảy mầm, hạt ngô không nảy mầm  là $0{,}28$.
    \itemch Xác suất hạt đậu không nảy mầm, hạt ngô nảy mầm là $0{,}18$.
  \end{itemchoice}
  }
\end{ex}
